% ---
% id: receiver/rem:e3
% kind: rem
% title: "Order dependence without support comparability"
% summary: >-
%   An acyclic graph of a rotation and a polarizer on the same two modes, with
%   no edge between them, compiles at generic angles to two different
%   contractions under its two sorts. Acyclicity alone therefore does not make a declared graph determine
%   an instrument.
% prominence: standard
% status: accepted
% uses: []
% assumes: []
% refs: []
% provenance: newly-derived
% lean:
% ---

\begin{remark}
\label{rem:e3}
Take $\Modes=\{m\}$ and two elements $e,f$ on $\Sigma_e=\Sigma_f=\{m\}$ with no
edge between them, $\Gmat_{\Sigma_e}=(\tfrac12)$ and
$\Gmat_{\Sigma_f}=(\tfrac12)$: here the two orders agree, and the example
looks harmless. Now take $\Modes=\{m_1,m_2\}$ with
$\Sigma_e=\Sigma_f=\Modes$, $\Gmat_{\Sigma_e}=\Rot(\theta)$ a rotation and
$\Gmat_{\Sigma_f}=\operatorname{diag}(1,0)$ a polarizer. The graph
$\Grph=(\{e,f\},\varnothing)$ is acyclic and every element is admissible, yet
$\Gmat_{\Sigma_f}\Gmat_{\Sigma_e}\neq\Gmat_{\Sigma_e}\Gmat_{\Sigma_f}$ for
$\theta\notin\frac{\pi}{2}\mathbb{Z}$: rotating and then projecting is not
projecting and then rotating. Two topological sorts of the \emph{same} declared
receiver would compile to two different instruments, and, because both are
contractions with the same declared port map, both would pass every other
validation in this contract. Without (E3), ``the declared receiver'' does not
denote.
\end{remark}
